\begin{longtable}{lcccccccc}
\toprule
Sismo & Fecha & $M$ & $d$ (km) & \multicolumn{2}{c}{SNR (dB)} & \multicolumn{2}{c}{SNR 10 s de P (dB)} & $T_{\text{mar}}$ (s)\\
\cmidrule(lr){5-6}\cmidrule(lr){7-8}
 & & & & antes & después & antes & después & \\
\midrule\endhead
Chiapas 1993 & 1993-09-10 & 7.2 & 863 & 39.8 & 50.0 & 13.3 & 22.9 & 9.6\\
Copala 1995 & 1995-09-14 & 7.4 & 289 & 61.5 & 76.4 & 43.1 & 58.1 & 3.9\\
Chiapas 1995 & 1995-10-21 & 7.2 & 664 & 50.6 & 61.7 & 35.6 & 46.6 & 5.6\\
Michoacán 1997 & 1997-01-11 & 7.2 & 397 & 66.3 & 74.5 & 41.2 & 49.4 & 4.4\\
Tehuacán 1999 & 1999-06-15 & 7.0 & 211 & 65.4 & 75.5 & 54.6 & 64.7 & 7.5\\
Oaxaca 1999 & 1999-09-30 & 7.5 & 433 & 47.8 & 58.8 & 36.5 & 47.6 & 3.5\\
Tecomán 2003 & 2003-01-22 & 7.6 & 522 & 67.0 & 77.2 & 39.5 & 49.7 & 6.9\\
El Mayor-Cucapah 2010 & 2010-04-04 & 7.2 & 2\,156 & 35.7 & 45.7 & 9.8 & 19.2 & 6.0\\
Ometepec 2012 & 2012-03-20 & 7.4 & 330 & 62.8 & 80.3 & 46.0 & 63.4 & 7.1\\
Golfo de California 2012 & 2012-04-12 & 7.0 & 1\,754 & 8.2 & 16.7 & -2.1 & -1.5 & 9.6\\
Papanoa 2014 & 2014-04-18 & 7.2 & 286 & 57.2 & 60.5 & 38.7 & 42.0 & 3.0\\
Chiapas 2014 & 2014-07-07 & 6.9 & 878 & 55.0 & 63.7 & 32.6 & 41.3 & 7.7\\
Tehuantepec 2017 & 2017-09-08 & 8.2 & 736 & 75.1 & 81.5 & 39.0 & 45.3 & 3.7\\
Puebla-Morelos 2017 & 2017-09-19 & 7.1 & 113 & 59.8 & 77.4 & 59.7 & 77.3 & 5.6\\
Pinotepa 2018 & 2018-02-16 & 7.2 & 350 & 64.5 & 76.7 & 39.6 & 51.8 & 6.7\\
Huatulco 2020 & 2020-06-23 & 7.4 & 508 & 68.0 & 81.4 & 42.4 & 55.7 & 7.5\\
Acapulco 2021 & 2021-09-08 & 7.0 & 271 & 72.8 & 80.3 & 59.2 & 66.7 & 7.3\\
Michoacán 2022 & 2022-09-19 & 7.6 & 410 & 68.2 & 78.9 & 28.7 & 39.4 & 8.0\\
Réplica Michoacán 2022 & 2022-09-22 & 6.8 & 415 & 71.9 & 80.1 & 45.9 & 54.2 & 6.2\\
Chiapas 2026 & 2026-07-17 & 7.3 & 846 & 45.6 & 55.2 & 22.4 & 31.9 & 6.0\\
\midrule
Promedio & & & & 57.2 & 67.6 & 36.3 & 46.3 & 6.4\textsuperscript{a}\\
\bottomrule
\caption{Resultados del método en los 20 sismos, registrados en la estación G.UNM. $d$ es la distancia epicentral y $T_{\text{mar}}$ el periodo del pico de microsismos del ruido. \textsuperscript{a}Mediana.}
\label{tab:sismos}
\end{longtable}